% =====================================================================
%  Chapitre 5 - Sprint 2 : referentiel de flotte et d'equipements
% =====================================================================
\chapter{Sprint 2 : referentiel de flotte et d'equipements}
\label{chap:sprint2}

\section*{Introduction}
\addcontentsline{toc}{section}{Introduction}

Le sprint precedent a etabli qui utilise la plateforme. Celui-ci etablit
\essentiel{avec quoi l'entreprise travaille} : le parc de vehicules, les
equipements embarques, les pieces administratives et le carburant. C'est
le second socle du systeme, celui sur lequel s'appuieront les missions et
la maintenance : on n'affecte pas un vehicule qui n'existe pas au
referentiel.

% =====================================================================
\section{Planification du sprint}

\subsection{Objectifs}

\begin{enumerate}
  \item Constituer un referentiel unique des vehicules et des
    equipements, chacun decrit une seule fois.
  \item Rattacher a chaque vehicule ses pieces administratives et ses
    assurances, avec leurs dates d'echeance.
  \item Suivre le cycle de vie d'un vehicule par son statut
    d'exploitation.
  \item Tracer l'affectation et la restitution des equipements.
  \item Enregistrer les pleins de carburant et en tirer une consommation.
\end{enumerate}

\subsection{Carnet du sprint}

\begin{table}[H]
\centering
\caption{Carnet du sprint 2}
\label{tab:backlog-s2}
\small
\begin{tabularx}{\textwidth}{@{}cp{2.5cm}Xcc@{}}
\toprule
\textbf{N\textsuperscript{o}} & \textbf{Acteur} & \textbf{Besoin exprime} &
\textbf{Prio.} & \textbf{Pts} \\ \midrule
1 & Utilisateur & Consulter le parc et rechercher un vehicule disponible & Must & 5 \\
2 & Responsable & Tenir la fiche d'un vehicule & Must & 8 \\
3 & Responsable & Enregistrer les pieces administratives et les assurances & Must & 5 \\
4 & Utilisateur & Consulter les echeances a venir & Must & 3 \\
5 & Responsable & Tenir la fiche d'un equipement & Must & 5 \\
6 & Utilisateur & Affecter et restituer un equipement & Must & 5 \\
7 & Utilisateur & Realiser une inspection d'equipement & Should & 3 \\
8 & Utilisateur & Declarer un plein de carburant & Should & 3 \\
9 & Utilisateur & Suivre la consommation de la flotte & Should & 5 \\
10 & Administrateur & Sortir un bien de l'inventaire & Must & 3 \\
\midrule
\multicolumn{4}{@{}r}{\textbf{Total}} & \textbf{45} \\
\bottomrule
\end{tabularx}
\end{table}

% =====================================================================
\section{Elaboration technique}

\subsection{Modelisation fonctionnelle}

\figReelle{fleet-usecase.png}
  {Cas d'utilisation du service de gestion de flotte}
  {fig:uc-fleet}{0.75}

La figure~\ref{fig:uc-fleet} distingue trois domaines. \essentiel{L'exploitation
courante} rassemble ce qui est ouvert a tout utilisateur authentifie : la
consultation du parc ne porte en effet aucune restriction de role dans le
code, seules les ecritures en sont assorties. \essentiel{La tenue du
referentiel} revient au responsable, et \essentiel{le retrait definitif}
au seul administrateur.

Cette gradation reproduit une regle simple : \essentiel{lire est ouvert,
ecrire est encadre, supprimer est reserve}. Elle se retrouvera dans tous
les modules suivants.

\begin{table}[H]
\centering
\caption{Description textuelle du cas « immatriculer un vehicule »}
\label{tab:desc-immat}
\small
\begin{tabularx}{\textwidth}{@{}p{2.5cm}X@{}}
\toprule
\textbf{Rubrique} & \textbf{Contenu} \\ \midrule
Acteur & Responsable \\
Precondition & Le vehicule n'existe pas encore au referentiel \\
Scenario nominal & Saisie de l'immatriculation, du numero de chassis, de
  la marque, du modele et des caracteristiques techniques ; controle
  d'unicite de l'immatriculation et du chassis ; enregistrement au statut
  DISPONIBLE \\
Alternatives & Immatriculation deja presente : creation refusee. Champ
  obligatoire manquant : le formulaire signale la ligne fautive \\
Postcondition & Le vehicule est visible au parc et peut recevoir des
  equipements, des documents et des missions \\
\bottomrule
\end{tabularx}
\end{table}

\subsection{Conception statique}

\figReelle{fleet-class-domain.png}
  {Modele du domaine du service de gestion de flotte}
  {fig:dom-fleet}{0.90}

Le modele s'organise autour de la classe \exemple{Vehicule}, a laquelle
se rattachent quatre familles d'objets : les \exemple{documents
administratifs} et les \exemple{assurances}, qui portent des dates
d'echeance ; les \exemple{equipements}, affectes puis restitues ; les
\exemple{pleins de carburant}, rattaches a une source ; et les
\exemple{inspections}.

Deux points de conception meritent d'etre releves.

\paragraph{L'echeance est une donnee, pas un calcul.}
La date de fin de validite est stockee sur le document et sur
l'assurance. Le service d'alertes la lit pour produire ses avertissements
plutot que de la deduire : une echeance depassee reste visible meme si
personne n'a ouvert l'ecran depuis des semaines.

\paragraph{L'equipement a son propre cycle de vie.}
Il porte un etat, un compteur d'utilisations et des seuils d'usure. Il
peut etre affecte a un vehicule ou rester au magasin : la relation est
donc optionnelle des deux cotes, ce que traduit la multiplicite du
diagramme.

\subsection{Conception dynamique}

% Ces deux sequences occupaient deja la meme largeur, la justification
% entiere. Elles paraissaient pourtant de tailles tres differentes, et
% elles l'etaient : la premiere mesure 1475 pixels de large, la seconde
% 895. A largeur egale, la premiere subit donc une reduction deux fois
% plus forte, et son texte devient deux fois plus petit.
%
% Egaliser les largeurs ne pouvait rien y faire. C'est le RAPPORT DE
% RENDU qu'il faut egaliser. La premiere passe en paysage : 24,7 cm au
% lieu de 16, soit un rendu a 0,63 de sa taille native. La seconde,
% laissee en pleine justification, rend a 0,68. Les deux textes ont
% desormais la meme taille a un vingtieme pres.
\figPaysage{fleet-seq-immatriculation.png}
  {Sequence : immatriculation d'un vehicule}
  {fig:seq-immat}

\figReelle{fleet-seq-suivi-echeances.png}
  {Sequence : consultation des echeances a venir}
  {fig:seq-echeances}{1.00}

La premiere sequence montre le controle d'unicite effectue avant tout
enregistrement. La seconde, figure~\ref{fig:seq-echeances}, illustre un
mecanisme different : \essentiel{aucune ecriture n'a lieu}, le service
agrege les dates portees par les documents, les assurances et les permis,
puis les ordonne par proximite. C'est une lecture transversale, et c'est
elle qui alimentera le module d'alertes du sprint~5.

Une troisieme sequence, le changement d'etat d'un vehicule, suit la meme
structure que l'immatriculation et n'est pas reproduite ici.

% =====================================================================
\section{Realisation}

\subsection{Le parc et l'inventaire}

La figure~\ref{fig:s2-vehicules} presente l'ecran de tete du module : le
parc, dans son ensemble. Chaque ligne porte l'immatriculation, la marque,
le modele, le kilometrage et surtout le \essentiel{statut
d'exploitation} --- disponible, en mission, en maintenance, hors service.
Ce statut est la donnee la plus consultee de tout le referentiel, parce
que c'est elle qui conditionne l'affectation : un vehicule en maintenance
ne doit pas apparaitre dans la liste des candidats a une mission. Les
filtres places en tete permettent de restreindre l'affichage a un statut,
a un type de vehicule ou a une plage de kilometrage.

\capture{s2-liste-vehicules.jpg}
  {Parc de vehicules : liste filtrable et statuts d'exploitation}
  {fig:s2-vehicules}

L'inventaire des equipements, figure~\ref{fig:s2-equipements}, obeit a la
meme logique de presentation mais repond a une question differente :
\exemple{ou se trouve ce materiel ?} La colonne d'affectation indique soit
un vehicule, soit le magasin. Un hayon, un groupe frigorifique ou un jeu
de sangles passe de l'un a l'autre au fil des besoins, et c'est ce
mouvement que l'ecran doit rendre lisible. Le compteur d'utilisations,
visible sur la droite, prepare la maintenance preventive du sprint~4 : au
dela d'un seuil, l'equipement doit etre controle.

\capture{s2-liste-equipements.jpg}
  {Inventaire des equipements et leur affectation courante}
  {fig:s2-equipements}

\subsection{La fiche d'un vehicule}

La fiche detaillee est organisee en onglets, pour eviter un formulaire
unique de plusieurs ecrans de haut. Le premier, figure~\ref{fig:s2-fiche1},
rassemble l'identite technique : chassis, energie, charge utile,
kilometrage, date de mise en circulation. Ces valeurs changent peu, mais
elles alimentent des decisions frequentes --- la charge utile determine
quelles marchandises le vehicule peut prendre, l'energie sert au calcul de
consommation.

\capture{s2-fiche-vehicule-1.jpg}
  {Fiche d'un vehicule : identite et caracteristiques techniques}
  {fig:s2-fiche1}

Le second onglet, figure~\ref{fig:s2-fiche2}, porte les pieces
administratives : carte grise, visite technique, vignette, autorisations
de transport. Chacune est datee et possede une \essentiel{fin de
validite}. C'est ici que se materialise le choix de conception expose plus
haut : la date est saisie une fois, stockee, puis relue par le module
d'alertes. L'ecran signale par une pastille de couleur les pieces dont
l'echeance approche, ce qui donne au responsable une lecture immediate
sans qu'il ait a comparer les dates lui-meme.

\capture{s2-fiche-vehicule-2.jpg}
  {Fiche d'un vehicule : pieces administratives et echeances}
  {fig:s2-fiche2}

Enfin, la figure~\ref{fig:s2-detail-equip} montre le detail d'un
equipement affecte. On y retrouve son etat, son compteur d'utilisations,
le vehicule porteur et l'historique de ses affectations successives. Cet
historique n'est pas un ornement : lorsqu'un equipement tombe en panne, il
permet de savoir sur quels vehicules il a servi et donc quelles missions
ont pu en patir.

\capture{s2-detail-equipement.jpg}
  {Detail d'un equipement et historique de ses affectations}
  {fig:s2-detail-equip}

% =====================================================================
\section{Tests de validation}

\begin{table}[H]
\centering
\caption{Tests de validation du sprint 2}
\label{tab:tests-s2}
\small
\begin{tabularx}{\textwidth}{@{}cXp{3.0cm}c@{}}
\toprule
\textbf{N\textsuperscript{o}} & \textbf{Cas teste} & \textbf{Attendu} &
\textbf{Etat} \\ \midrule
T1 & Immatriculation d'un vehicule inedit & Vehicule cree, statut DISPONIBLE & Reussi \\
T2 & Immatriculation deja presente & Creation refusee & Reussi \\
T3 & Document dont l'echeance est depassee & Vehicule signale dans les echeances & Reussi \\
T4 & Affectation d'un equipement a un vehicule & Equipement indisponible au magasin & Reussi \\
T5 & Restitution d'un equipement & Compteur d'utilisations incremente & Reussi \\
T6 & Plein de carburant enregistre & Consommation moyenne recalculee & Reussi \\
\bottomrule
\end{tabularx}
\end{table}

% =====================================================================
\section*{Conclusion}
\addcontentsline{toc}{section}{Conclusion}

Ce sprint livre le referentiel materiel de la plateforme : seize
vehicules decrits avec leurs caracteristiques, leurs pieces
administratives et leurs assurances, quatorze equipements suivis avec
leur etat et leur affectation, et le suivi du carburant. Les
quarante-cinq points du carnet ont ete realises.

L'apport le plus durable de ce sprint n'est pas un ecran mais une
donnee : \essentiel{les dates d'echeance sont desormais stockees et
interrogeables}. Sans elles, la surveillance reglementaire du sprint~5
n'aurait aucune matiere sur laquelle travailler.

Les deux socles etant poses, le sprint suivant aborde le coeur de
l'activite : la chaine qui va de la demande du client a la livraison de
sa marchandise.
